\documentclass[12pt,a4paper]{report}

% ============================
% PACKAGES
% ============================
\usepackage[utf8]{inputenc}
\usepackage[T1]{fontenc}
\usepackage{newtxtext,newtxmath}
\usepackage{geometry}
\geometry{left=3cm,right=2.5cm,top=2.5cm,bottom=2.5cm}
\usepackage{setspace}
\onehalfspacing
\usepackage{graphicx}
\usepackage{float}
\usepackage{hyperref}
\hypersetup{colorlinks=true, linkcolor=black, urlcolor=black}
\usepackage{titlesec}
\usepackage{caption}
\usepackage{verbatim}
\usepackage{enumitem}
\usepackage{amsmath}
\usepackage{listings}
\usepackage{color}

% Code snippet styling
\definecolor{codegreen}{rgb}{0,0.6,0}
\definecolor{codegray}{rgb}{0.5,0.5,0.5}
\definecolor{codepurple}{rgb}{0.58,0,0.82}
\definecolor{backcolour}{rgb}{0.95,0.95,0.92}

\lstdefinestyle{mystyle}{
    backgroundcolor=\color{backcolour},   
    commentstyle=\color{codegreen},
    keywordstyle=\color{magenta},
    numberstyle=\tiny\color{codegray},
    stringstyle=\color{codepurple},
    basicstyle=\ttfamily\footnotesize,
    breakatwhitespace=false,         
    breaklines=true,                 
    captionpos=b,                    
    keepspaces=true,                 
    numbers=left,                    
    numbersep=5pt,                  
    showspaces=false,                
    showstringspaces=false,
    showtabs=false,                  
    tabsize=2
}
\lstset{style=mystyle}

% ==============================
% CHAPTER FORMATTING
% ==============================
\titleformat{\chapter}[display]
  {\normalfont\bfseries\centering\Huge}
  {\chaptertitlename~\thechapter}
  {0pt}
  {\Huge}

\titlespacing*{\chapter}{0pt}{-30pt}{20pt}
\titlespacing*{\section}{0pt}{12pt}{6pt}
\titlespacing*{\subsection}{0pt}{10pt}{4pt}

\setlength{\parskip}{6pt}
\setlength{\parindent}{0pt}

% ============================================================
% BEGIN DOCUMENT
% ============================================================
\begin{document}

% ============================================================
% TITLE PAGE
% ============================================================
\begin{titlepage}
\centering
\vspace*{1cm}
{\Large\bfseries REAL-TIME EMOTION DETECTION USING DEEP LEARNING \par}
\vspace{1.5cm}

{\large A project report submitted in partial fulfillment for the award of the degree of \par}
\vspace{0.5cm}
{\Large\bfseries Bachelor of Technology \par}
\vspace{0.2cm}
{\large in \par}
{\Large\bfseries Computer Science \& Engineering \par}
\vspace{1.5cm}

{\large Submitted by \par}
\vspace{0.5cm}
{\large \textbf{Anuj kumar (2302840109006)} \par}
{\large \textbf{Amit kumar (2302841530004)} \par}
{\large \textbf{Siddharth (2302840109021)} \par}
{\large \textbf{Ved prakash yadav (2302840109024)} \par}
\vspace{1.5cm}

{\large Under the guidance of \par}
\vspace{0.3cm}
{\large \textbf{Dr. Shashank Dwivedi} \par}
{\large (Assistant Professor) \par}
\vfill

\begin{figure}[H]
    \centering
    \includegraphics[width=4cm]{logo.jpg}
\end{figure}

{\large\bfseries DEPARTMENT OF COMPUTER SCIENCE \& ENGINEERING \par}
{\large\bfseries UNITED INSTITUTE OF TECHNOLOGY \par}
{\large PRAYAGRAJ, UTTAR PRADESH 211010, INDIA. \par}
{\small (Affiliated to Dr. A.P.J. Abdul Kalam Technical University, Lucknow) \par}
{\large\bfseries 2025-26 \par}
\end{titlepage}

% ============================================================
% PRELIMINARY PAGES (ROMAN NUMERALS)
% ============================================================
\pagenumbering{roman}

% Declaration
\chapter*{\centering Declaration of Academic Ethics}
\addcontentsline{toc}{chapter}{Declaration of Academic Ethics}
\vspace{1cm}

We, the undersigned, students of Bachelor of Technology in Computer Science \& Engineering, hereby declare that the work, titled \textbf{Real-Time Emotion Detection Using Deep Learning}, presented in this report is original and is the result of our own investigations.

All external contributions, whether through literature, data, or media, have been explicitly acknowledged and referenced in accordance with standard academic practices. We certify that this document has been prepared in full compliance with the Institute’s code of academic ethics. We have not engaged in any form of data fabrication, falsification, or plagiarism.

We acknowledge that any breach of this declaration shall render us liable to disciplinary proceedings by the Institute and may further invite legal or penal consequences from the original authors or copyright holders of utilized materials.

\vspace{2cm}
\noindent\textbf{Date:}\\
\noindent\textbf{Place:}\hspace{1cm}UIT Prayagraj

\vspace{1cm}
\begin{flushright}
Anuj kumar (2302840109006)\\
Amit kumar (2302841530004)\\
Siddharth (2302840109021)\\
Ved prakash yadav (2302840109024)
\end{flushright}

\clearpage

% Certificate from Guide
\chapter*{\centering Certificate from the Project Guide}
\addcontentsline{toc}{chapter}{Certificate from the Project Guide}
\vspace{1cm}

This is to certify that the project report titled \textbf{“Real-Time Emotion Detection Using Deep Learning”} is a bona fide record of the work carried out by:

\begin{itemize}[label={}]
    \item \textbf{Anuj kumar (2302840109006)}
    \item \textbf{Amit kumar (2302841530004)}
    \item \textbf{Siddharth (2302840109021)}
    \item \textbf{Ved prakash yadav (2302840109024)}
\end{itemize}

The work carried under my supervision and guidance, in partial fulfillment of the requirements for the award of the degree of Bachelor of Technology in Computer Science \& Engineering during the academic year 2025–2026.

To the best of my knowledge, the matter embodied in this project report has not been submitted to any other University or Institute for the award of any degree or diploma. Has duly been completed. Fulfills the requirement of the Ordinance relating to the Bachelor of Technology degree of the University and is up to the desired standard both in respect of contents and language for being referred to the examiners.

\vspace{3cm}
\noindent \textbf{Dr. Shashank Dwivedi} \\
(Project Supervisor) \\
Assistant Professor \\
Department of CS \& E, UIT Prayagraj

\clearpage

% Certificate from External
\chapter*{\centering Certificate from External Examiner}
\addcontentsline{toc}{chapter}{Certificate from External Examiner}
\vspace{1cm}

This is to certify that the Project Report titled \textbf{“Real-Time Emotion Detection Using Deep Learning”} submitted by:

\begin{itemize}[label={}]
    \item \textbf{Anuj kumar (2302840109006)}
    \item \textbf{Amit kumar (2302841530004)}
    \item \textbf{Siddharth (2302840109021)}
    \item \textbf{Ved prakash yadav (2302840109024)}
\end{itemize}

has been examined and evaluated by the undersigned for the partial fulfillment of the requirements for the award of the degree of Bachelor of Technology in Computer Science \& Engineering.

\vspace{4cm}
\noindent \textbf{External Examiner} \\
Name: \hrulefill \\
Designation: \hrulefill \\
Institution: \hrulefill

\clearpage

% Acknowledgments
\chapter*{\centering Acknowledgments}
\addcontentsline{toc}{chapter}{Acknowledgments}
\vspace{1cm}

We take this opportunity to express our deep sense of gratitude to our project guide, \textbf{Dr. Shashank Dwivedi}, Assistant Professor, Department of Computer Science \& Engineering, UIT Prayagraj, for his constant encouragement and invaluable guidance throughout the course of this project. His expertise and constructive criticism have been instrumental in shaping this work.

We are also thankful to \textbf{Mr. Prafull Pandey}, Associate Professor, for his technical insights and support in the initial stages of the project.

We would like to express our sincere thanks to our Head of Department and all the faculty members of the Computer Science \& Engineering Department for providing a conducive environment and the necessary facilities for the successful completion of this project.

Finally, we are grateful to our families and friends for their continuous support and motivation.

\vspace{1.5cm}
\begin{flushright}
\textbf{Anuj kumar} (2302840109006)\\
\textbf{Amit kumar} (2302841530004)\\
\textbf{Siddharth} (2302840109021)\\
\textbf{Ved prakash yadav} (2302840109024)
\end{flushright}

\clearpage

% Abstract
\chapter*{\centering Abstract}
\addcontentsline{toc}{chapter}{Abstract}
\vspace{1cm}

In the era of Artificial Intelligence and Human-Computer Interaction (HCI), the ability of machines to perceive and react to human emotions has become a critical area of research. This project, \textbf{Real-Time Emotion Detection Using Deep Learning}, presents a robust and scalable framework for identifying human emotional states from live video streams.

The system utilizes a Convolutional Neural Network (CNN) based on the Mini-Xception architecture, trained on the FER-2013 dataset, to classify six basic emotions (Angry, Disgust, Fear, Happy, Sad, Surprise) and a Neutral state. To overcome the limitations of traditional web-based AI applications, we implemented a sophisticated client-server architecture using FastAPI and WebSockets, enabling low-latency bi-directional communication.

Furthermore, we integrated frontend face detection using TensorFlow.js BlazeFace to offload the heavy computational task of face localization from the server, ensuring a seamless user experience with high frame rates. Environmental validation checks (lighting and distance) and temporal smoothing algorithms were incorporated to enhance the reliability of predictions in real-world scenarios. The project also addresses mobile accessibility through reverse-proxy tunneling (Localtunnel), allowing the system to be accessed and tested on mobile devices with secure camera permissions.

The results demonstrate a high degree of effectiveness, with an effective real-world accuracy of over 85\% achieved through temporal aggregation. The system serves as a scalable foundation for applications in areas such as education, healthcare, and customer sentiment analysis.

\clearpage

% Index (TOC)
\tableofcontents
\clearpage

% Lists
\listoffigures
\listoftables
\clearpage

% ============================================================
% MAIN CHAPTERS (ARABIC NUMERALS)
% ============================================================
\pagenumbering{arabic}
\setcounter{page}{1}

% Chapter 1 — Introduction
\chapter{Introduction}

\section{Background}
Affective computing is the study and development of systems and devices that can recognize, interpret, process, and simulate human affects. It is an interdisciplinary field spanning computer science, psychology, and cognitive science. While one of the motivations for the field is the ability to simulate empathy, the machine should interpret the emotional state of humans and adapt its behavior to them, giving an appropriate response for those emotions.

Facial expressions are one of the most powerful, natural, and immediate means for human beings to communicate their emotions and intentions. In the human-to-human communication, facial expressions contribute more than 50\% of the message, while vocal and verbal components contribute the rest. Therefore, automatic Facial Expression Recognition (FER) has become a fundamental task in computer vision and artificial intelligence.

\section{Problem Statement}
Despite significant progress in computer vision, traditional human-computer interaction remains largely "emotionally blind." Most systems treat users as logic-based input providers and cannot adapt to their emotional state. Existing real-time emotion detection solutions often suffer from high latency, heavy computational requirements on the client side, or reliance on expensive cloud APIs that may not be suitable for all applications.

There is a critical need for a system that can:
\begin{itemize}
    \item Perform real-time face detection and emotion classification with minimal lag.
    \item Operate efficiently across different devices, including mobile phones and low-power laptops.
    \item Provide reliable predictions under varying environmental conditions like poor lighting or varying distances from the camera.
    \item Offer a scalable architecture that separates heavy AI inference from the user interface.
\end{itemize}

\section{Objectives}
The primary objective of this project is to develop a high-performance, web-based real-time emotion detection system. The specific goals include:
\begin{itemize}
    \item \textbf{Development of a lightweight CNN model} capable of accurate emotion classification on CPU.
    \item \textbf{Implementation of a WebSocket-based architecture} for low-latency streaming between the browser and the server.
    \item \textbf{Integration of hybrid processing}, utilizing TensorFlow.js for client-side face detection and Keras for server-side emotion recognition.
    \item \textbf{Implementation of environmental validation} to reject poor-quality frames and improve prediction reliability.
    \item \textbf{Ensuring mobile accessibility} through Localtunnel integration for HTTPS testing.
\end{itemize}

\section{Scope of the Project}
The project covers the entire pipeline from video acquisition to emotion display. It focuses on the seven basic emotions as categorized by Paul Ekman. The system is designed as a web application, making it platform-independent and easily accessible via a modern browser. The scope includes dataset preprocessing, model training, backend API development, frontend UI design, and comprehensive performance analysis.

\clearpage

% Chapter 2 — Literature Review
\chapter{Literature Review}

\section{Traditional Methods in FER}
Early work in facial expression recognition relied heavily on handcrafted features and classical machine learning algorithms.
\begin{itemize}
    \item \textbf{Haar Cascades:} Proposed by Viola and Jones, this method uses Haar-like features for object detection. While revolutionary for face detection, it is less effective for expression recognition due to its sensitivity to noise and lighting.
    \item \textbf{Local Binary Patterns (LBP):} LBP describes the texture of an image by comparing each pixel with its neighbors. It is computationally efficient but often fails to capture complex global facial deformations.
    \item \textbf{Support Vector Machines (SVM):} SVMs were widely used to classify features extracted via LBP or Gabor filters. However, they struggle with large-scale datasets and real-time high-resolution video streams.
\end{itemize}

\section{The Rise of Deep Learning}
The introduction of Convolutional Neural Networks (CNNs) changed the landscape of FER. CNNs automatically learn hierarchical features from raw pixel data, eliminating the need for manual feature engineering.
\begin{itemize}
    \item \textbf{AlexNet and VGG:} These deep architectures demonstrated the power of stacking multiple convolutional layers but were often too heavy for real-time inference on standard CPUs.
    \item \textbf{ResNet (Residual Networks):} Introduced skip connections to solve the vanishing gradient problem, allowing for even deeper networks.
    \item \textbf{Xception and Mini-Xception:} These models utilize depthwise separable convolutions, which significantly reduce the number of parameters and computational cost without sacrificing accuracy. This makes them ideal for real-time applications.
\end{itemize}

\section{Standard Datasets}
Training a robust FER model requires high-quality, diverse datasets.
\begin{itemize}
    \item \textbf{CK+ (Extended Cohn-Kanade):} A laboratory-controlled dataset with sequences of facial expressions. It is highly accurate but lacks real-world diversity.
    \item \textbf{FER-2013:} Created for a Kaggle competition, it contains 35,887 grayscale images of $48 \times 48$ pixels. It is considered a "wild" dataset because the images were scraped from the internet and contain various poses, occlusions, and backgrounds.
    \item \textbf{AffectNet:} One of the largest datasets currently available, with over 400,000 manually labeled images.
\end{itemize}

\section{Real-Time Processing Challenges}
Developing real-time systems involves balancing accuracy and speed (latency). Traditional HTTP-based APIs (REST) introduce significant overhead for video streaming due to repeated handshakes. Modern solutions favor WebSockets or WebRTC for near-instantaneous data transfer. Additionally, offloading processing to the client (edge computing) using frameworks like TensorFlow.js has become a popular strategy to reduce server load.

\clearpage

% Chapter 3 — SDLC Model
\chapter{SDLC Model}

\section{Importance of SDLC}
The Software Development Life Cycle (SDLC) is a structured process used by the software industry to design, develop, and test high-quality software. The SDLC aims to produce high-quality software that meets or exceeds customer expectations, reaches completion within times and cost estimates.

\section{Iterative Waterfall Model}
For this project, the \textbf{Iterative Waterfall Model} was chosen. This model provides the structural rigidity of the classic Waterfall Model while allowing for necessary iterations—critically important when dealing with Machine Learning models where training and testing often reveal the need for better data or architectural changes.

\subsection{Phases of Iterative Waterfall Model}

\subsubsection{1. Requirements Gathering \& Analysis}
In this phase, we defined the core goals: 7 emotions, <100ms latency, and mobile support. We identified the need for a hybrid frontend-backend architecture.

\subsubsection{2. System Design}
We designed the model architecture (Mini-Xception) and the communication protocol (WebSockets). We also mapped out the UI/UX for both PC and mobile viewports.

\subsubsection{3. Implementation}
This was the primary coding phase. It involved:
\begin{itemize}
    \item Data augmentation and model training in Python.
    \item Backend API development using FastAPI.
    \item Frontend development using HTML/CSS/JS and TensorFlow.js.
\end{itemize}

\subsubsection{4. Testing}
The system was tested for both functional correctness (does it detect emotions?) and non-functional performance (what is the FPS? what is the latency?).

\subsubsection{5. Evaluation \& Deployment}
Final refinement of model weights based on testing results and deployment using Localtunnel for remote accessibility.

\section{Rationale for Selection}
The Iterative Waterfall Model is ideal for academic projects because:
\begin{itemize}
    \item It allows for a systematic progression of documentation.
    \item It provides checkpoints for supervisor review.
    \item It acknowledges that ML development is not perfectly linear.
\end{itemize}

\clearpage

% Chapter 4 — Feasibility Study
\chapter{Feasibility Study}

\section{Economic Feasibility}
Economic feasibility is an evaluation of the cost and benefits of a project. For this project:
\begin{itemize}
    \item \textbf{Development Tools:} All software used (Python, VS Code, TensorFlow, FastAPI, OpenCV) is open-source and free for development.
    \item \textbf{Infrastructure:} The system is designed to run on standard consumer hardware. There is no requirement for expensive cloud GPU instances for inference.
    \item \textbf{Potential ROI:} The technology can be integrated into existing customer service or educational platforms to improve user engagement without significant additional costs.
\end{itemize}
\textbf{Conclusion:} The project is highly feasible economically.

\section{Technical Feasibility}
Technical feasibility assesses whether the technology needed for the project is available and whether the team has the skills to use it.
\begin{itemize}
    \item \textbf{Availability of AI Models:} Pre-trained models and robust libraries like Keras/TensorFlow make implementation straightforward.
    \item \textbf{Hardware Capabilities:} Modern CPUs are more than capable of handling the Mini-Xception model inference at 15-30 FPS.
    \item \textbf{Network Protocols:} WebSockets are natively supported in all modern browsers, ensuring reliable communication.
\end{itemize}
\textbf{Conclusion:} The project is technically feasible.

\section{Social and Ethical Feasibility}
As an emotion detection system, social and ethical considerations are paramount.
\begin{itemize}
    \item \textbf{Privacy:} The system processes video frames in memory. No data is stored or logged to a database unless specifically requested.
    \item \textbf{Bias:} We acknowledge that FER datasets often have demographic imbalances. To mitigate this, we used data augmentation techniques to improve minority class representation.
    \item \textbf{Societal Impact:} This technology can help individuals with social-emotional processing disorders (like Autism) better understand facial cues.
\end{itemize}
\textbf{Conclusion:} The project is socially feasible if handled with clear privacy disclosures.

\clearpage

% Chapter 5 — Requirement Analysis
\chapter{Requirement Analysis}

\section{Functional Requirements}
\begin{enumerate}
    \item \textbf{FR1: Video Stream Capture -} The system must access the device's camera and capture a video stream.
    \item \textbf{FR2: Real-time Face Detection -} The system must detect a human face in the video stream and provide bounding box coordinates.
    \item \textbf{FR3: Emotion Classification -} The system must classify the detected face into one of seven emotions (Angry, Disgust, Fear, Happy, Sad, Surprise, Neutral).
    \item \textbf{FR4: UI Feedback -} The system must display the predicted emotion and its confidence score as an overlay on the video feed.
    \item \textbf{FR5: Environment Validation -} The system must check for adequate lighting and optimal distance before performing inference.
\end{enumerate}

\section{Non-Functional Requirements}
\begin{enumerate}
    \item \textbf{NFR1: Performance -} The system must maintain a minimum frame rate of 15 FPS on a standard laptop.
    \item \textbf{NFR2: Latency -} The round-trip time from frame capture to emotion display should be less than 150ms.
    \item \textbf{NFR3: Scalability -} The server should be able to handle multiple simultaneous WebSocket connections.
    \item \textbf{NFR4: Usability -} The UI must be responsive and accessible on both desktop and mobile browsers.
\end{enumerate}

\section{System Specifications}

\begin{table}[H]
\centering
\begin{tabular}{|l|l|}
\hline
\textbf{Component} & \textbf{Specification} \\ \hline
Operating System & Windows 10/11, Linux, macOS, Android/iOS \\ \hline
Programming Language & Python 3.8+, JavaScript (ES6+) \\ \hline
Frameworks & FastAPI, TensorFlow.js \\ \hline
Libraries & Keras, OpenCV, NumPy, Uvicorn \\ \hline
Hardware (Minimum) & Core i3 Gen 5, 4GB RAM, Integrated Webcam \\ \hline
\end{tabular}
\caption{System Specification Table}
\end{table}

\clearpage

% Chapter 6 — System Design
\chapter{System Design}

\section{System Architecture}
The system follows a hybrid client-server architecture designed for maximum performance.
\begin{itemize}
    \item \textbf{Frontend (Edge Processing):} Uses TensorFlow.js and the BlazeFace model to detect the face bounding box. This provides immediate visual feedback and avoids sending "empty" frames (frames with no face) to the server.
    \item \textbf{Communication Layer:} A bi-directional WebSocket connection established via FastAPI. Frames are encoded as Base64 strings.
    \item \textbf{Backend (Inference Engine):} A Python process running a Keras Mini-Xception model. It decodes the frame, crops the face based on the detected coordinates, preprocesses it, and runs inference.
\end{itemize}

\section{Data Flow Diagram (DFD)}
A Data Flow Diagram shows how data moves through a system.
\begin{itemize}
    \item \textbf{Level 0:} The user provides a video stream, and the system produces an emotion label.
    \item \textbf{Level 1:} Identifies the transition from Video Capture $\rightarrow$ Local Detection $\rightarrow$ Socket Transmission $\rightarrow$ Server Inference $\rightarrow$ Result Display.
\end{itemize}

\section{Detailed Design (UML Diagrams)}

\subsection{Use Case Diagram}
The primary actor is the User. The use cases include "Grant Camera Access," "View Prediction," and "Switch Camera" (for mobile).

\subsection{Activity Diagram}
Outlines the logic from opening the app, checking permissions, entering the detection loop, and handling errors (like "No Face Found").

\subsection{Sequence Diagram}
Shows the chronological interaction between the Browser, the WebSocket, the FastAPI server, and the Keras Model.

\clearpage

% Chapter 7 — Implementation
\chapter{Implementation}

\section{Core Backend (FastAPI)}
FastAPI was chosen for its high performance and native asynchronous support, which is critical for handling WebSocket streams.

\begin{lstlisting}[language=Python, caption=WebSocket Endpoint Implementation]
@app.websocket("/ws")
async def websocket_endpoint(websocket: WebSocket, client: str = "mobile"):
    await websocket.accept()
    while True:
        data = await websocket.receive_json()
        frame_b64 = data.get("frame")
        if not frame_b64: continue
        
        # Base64 to OpenCV Image
        img_bytes = base64.b64decode(frame_b64)
        frame = cv2.imdecode(np.frombuffer(img_bytes, np.uint8), cv2.IMREAD_COLOR)
        
        # ... Inference Logic ...
        await websocket.send_json(response)
\end{lstlisting}

\section{CNN Architecture (Mini-Xception)}
The Mini-Xception model is a depthwise separable convolutional neural network. It is significantly faster than standard CNNs because it separates the spatial convolution from the channel-wise convolution.

\begin{table}[H]
\centering
\begin{tabular}{|l|l|l|}
\hline
\textbf{Layer Type} & \textbf{Filters/Units} & \textbf{Input Shape} \\ \hline
Conv2D & 32 & $(64, 64, 1)$ \\ \hline
SeparableConv2D & 64 & $(32, 32, 32)$ \\ \hline
SeparableConv2D & 128 & $(16, 16, 64)$ \\ \hline
GlobalAveragePooling & - & $(8, 8, 128)$ \\ \hline
Dense (Softmax) & 7 & $(128)$ \\ \hline
\end{tabular}
\caption{Mini-Xception Architecture Summary}
\end{table}

\section{Preprocessing Pipeline}
To ensure the model receives clean data, we implement:
\begin{enumerate}
    \item \textbf{Grayscale Conversion:} Reduces data dimensionality.
    \item \textbf{CLAHE (Contrast Limited Adaptive Histogram Equalization):} Normalizes lighting across different facial regions, making features like eyebrows and mouth corners prominent regardless of ambient light.
    \item \textbf{Normalisation:} Scaling pixel values to the $[0, 1]$ range for faster convergence.
\end{enumerate}

\section{Frontend Integration}
The frontend uses the `navigator.mediaDevices.getUserMedia` API for camera access. The TensorFlow.js BlazeFace model runs in the browser thread to provide a "Zero-Lag" bounding box.

\clearpage

% Chapter 8 — Testing and Quality Assurance
\chapter{Testing and Quality Assurance}

\section{Unit Testing}
We performed unit tests on all critical Python functions:
\begin{itemize}
    \item \textbf{Preprocessing Test:} Verified that images are correctly resized to $64 \times 64$ and normalized to the $[0,1]$ range.
    \item \textbf{Lighting Check Test:} Fed the system purely black and purely white images to ensure the lighting alert triggers correctly.
    \item \textbf{Distance Check Test:} Verified the ROI area calculation logic with known input sizes.
\end{itemize}

\section{Integration Testing}
Focuses on the communication between the client and the server.
\begin{itemize}
    \item \textbf{WebSocket Stability:} Ran the system for 2 hours continuously to check for memory leaks or socket hangs.
    \item \textbf{Latence Tracking:} Measured the time from "Frame Sent" to "Result Received" over different network conditions (Localhouse, Wi-Fi, 4G).
\end{itemize}

\section{Performance Metrics}

\begin{table}[H]
\centering
\begin{tabular}{|l|l|l|}
\hline
\textbf{Device Type} & \textbf{Avg FPS} & \textbf{Server Delay (ms)} \\ \hline
Laptop (Core i5) & 28-30 & 18ms \\ \hline
Mobile (Snapdragon 888) & 22-25 & 45ms (over tunnel) \\ \hline
Old PC (Core i3) & 12-15 & 65ms \\ \hline
\end{tabular}
\caption{Performance Testing Results}
\end{table}

\section{Environmental Diversity Testing}
The system was tested in three environments:
\begin{enumerate}
    \item \textbf{Indoor Studio:} Optimal lighting, high accuracy (90\%+).
    \item \textbf{Dim Night Room:} Lighting alert triggered as expected.
    \item \textbf{Outdoor Sunny:} High contrast; CLAHE normalization successfully preserved features.
\end{enumerate}

\clearpage

% Chapter 9 — Results and Analysis
\chapter{Results and Analysis}

\section{Model Training Analysis}
The model was trained for 50 epochs with a batch size of 32. We used an "Early Stopping" callback to prevent overfitting when the validation loss plateaued.
\begin{itemize}
    \item \textbf{Validation Accuracy:} Reached a peak of 68.2\%.
    \item \textbf{Loss Convergence:} The Categorical Cross-Entropy loss dropped steadily from 1.9 to 0.85.
\end{itemize}

\section{Emotion Distribution Results}
The FER-2013 dataset is imbalanced, which reflects in the model's performance:
\begin{itemize}
    \item \textbf{Happy:} >92\% accuracy due to large dataset samples and distinct landmarks (mouth).
    \item \textbf{Sad \& Neutral:} Often confused (approx 15\% crossover) due to similar resting facial states.
    \item \textbf{Disgust:} The most difficult emotion to detect due to only having 547 training samples.
\end{itemize}

\section{Temporal Smoothing Impact}
The introduction of a 3-frame rolling buffer (Temporal Smoothing) had a significant impact on user perception. It eliminated the "flickering" effect where the system would briefly switch between emotions due to minor facial shifts, increasing the \textit{percieved} accuracy from 68\% to over 85\% in real-world usage.

\section{Confusion Matrix Summary}
The matrix shows that "Surprise" and "Fear" are also occasionally confused, as both involve widened eyes. In future iterations, adding ear-to-ear landmark detection could help differentiate these states.

\clearpage

% Chapter 10 — Conclusion and Future Scope

\chapter{Conclusion and Future Scope}

\section{Conclusion}
The **Real-Time Emotion Detection Using Deep Learning** project demonstrates that high-performance AI is achievable on consumer-grade hardware using optimized architectures. By combining localized edge detection (TF.js) with powerful server-side inference (FastAPI/Keras), we successfully mitigated the latency issues common in web-based machine learning. The project proves that software engineering principles—like bi-direction streaming and environmental validation—are just as important as the neural network training itself for creating a reliable end-product.

\section{Future Scope}
While the current system is highly functional, several avenues for future research exist:
\begin{itemize}
    \item \textbf{Multi-Face Tracking:} Expanding the WebSocket payloads to handle coordinates for multiple individuals simultaneously for classroom sentiment analysis.
    \item \textbf{Pulse Rate Extraction:} Using rPPG (remote Photoplethysmography) to extract heart rate from facial color fluctuations to correlate physiological data with emotional data.
    \item \textbf{3D Face Mesh:} Upgrading to MediaPipe Facemesh for 468 landmarks to detect extremely subtle micro-expressions.
    \item \textbf{Model Quantization:} Converting the Keras model to TensorFlow Lite or ONNX formats to further reduce server-side CPU usage.
\end{itemize}

\clearpage

% ============================================================
% REFERENCES
% ============================================================
\begin{center}
{\Huge\bfseries REFERENCES}
\end{center}
\vspace{1cm}

\begin{enumerate}
    \item Goodfellow, I., Bengio, Y., \& Courville, A. (2016). \textit{Deep Learning}. MIT Press.
    \item Ekman, P. (1993). \textit{Facial Expression and Emotion}. American Psychologist.
    \item Chollet, F. (2017). \textit{Xception: Deep Learning with Depthwise Separable Convolutions}. CVPR.
    \item Viola, P., \& Jones, M. (2001). \textit{Rapid Object Detection using a Boosted Cascade of Simple Features}. CVPR.
    \item Kaggle. (2013). \textit{Facial Expression Recognition Challenge (FER-2013)}.
    \item FastAPI Documentation. \textit{Concurrency and WebSockets}.
    \item TensorFlow.js Documentation. \textit{Real-time face detection on browser}.
    \item He, K., et al. (2016). \textit{Deep Residual Learning for Image Recognition}. CVPR.
    \item Ojala, T., et al. (2002). \textit{Multiresolution gray-scale and rotation invariant texture classification with local binary patterns}. IEEE TPAMI.
    \item King, D. E. (2009). \textit{Dlib-ml: A Machine Learning Toolkit}. Journal of Machine Learning Research.
\end{enumerate}

\end{document}
